\documentclass[12pt,a4paper]{article}
\usepackage[utf8]{inputenc}
\usepackage[T2A,T1]{fontenc}
\usepackage[ukrainian]{babel}
\usepackage{amsmath,amssymb}
\usepackage{booktabs}
\usepackage{longtable}
\usepackage{array}
\usepackage{calc}
\usepackage{geometry}
\geometry{margin=2.5cm}

\begin{document}
\section{РЕФЕРАТ}\label{ux440ux435ux444ux435ux440ux430ux442}



\textbf{Система та спосіб обробки багатоканальних EEG-даних для обчислення показника обчислювальної ментальної енергії (CME) із застосуванням квантового обчислювального модуля та метаевристичної оптимізації}



Винахід належить до галузі нейроінформатики, великих даних та медичних інформаційних технологій (МПК: A61B 5/372, G06N 10/00, G16H 50/20).



Запропоновано систему та спосіб обробки багатоканальних EEG-даних для обчислення стандартизованого показника обчислювальної ментальної енергії (CME). Система містить пристрій реєстрації EEG щонайменше з чотирма електродами, модулі сегментації, попередньої обробки та виділення спектральних ознак у діапазонах \(\delta/\theta/\alpha/\beta/\gamma\), квантовий обчислювальний модуль у вигляді 4-кубітної варіаційної квантової схеми з архітектурою повторного завантаження даних для оцінки ймовірності стану потоку \(p_{\text{flow}}(t)\), та модуль метаевристичної оптимізації параметрів квантової схеми. Показник CME обчислюється за формулою \(\text{CME}(t) = \kappa \cdot E_{\text{band}}(t) \cdot g(c, p_{\text{flow}}) \cdot \Delta\) та вимірюється в одиницях Вн (\(1\;\text{Вн} \equiv 1\;\text{мкВ}^2 \cdot \text{с}\)). Технічний результат полягає у формуванні стандартизованого вимірюваного показника когнітивного навантаження із керованістю ресурсів квантового процесора.



Кількість пунктів формули: 2 незалежних, 12 залежних. Кількість креслень: 4 (Фіг. 1--4). Фіг. 1.


\end{document}
